% context: markscheme (official solutions) for 2-1DN_Mental-math
% topic: mental math, no calculator, answers only — metric prefixes, fractions to
%        decimals, fractions of an hour and a day, time facts, inches to cm, km to miles
% convention: answers in bold; "=" joins exact equalities, "\approx" a rounded or
%        approximate value; italic notes give the one-line reasoning and the usual traps
\documentclass[12pt, twoside]{article}
%\documentclass[12pt, twoside]{article}
\usepackage[letterpaper, margin=1in, headsep=0.2in]{geometry}
\setlength{\headheight}{0.6in}
%\usepackage[english]{babel}
\usepackage[utf8]{inputenc}
\usepackage{microtype}
\usepackage{amsmath}
\usepackage{amssymb}
%\usepackage{amsfonts}
\usepackage[nomessages]{fp} %\FPeval{\var-name}{2*sin(pi/6)}
\usepackage{siunitx} %units in math. eg 20\milli\meter
\usepackage{yhmath} % for arcs, overparenth command
\usepackage{tikz} %graphics
\usetikzlibrary{quotes, angles, arrows, arrows.meta}
\usepackage{graphicx} %consider setting \graphicspath{{images/}}
\usepackage{parskip} %no paragraph indent
\usepackage{enumitem}
\usepackage{multicol}
\usepackage{venndiagram}

\usepackage{fancyhdr}
\pagestyle{fancy}
\fancyhf{}
\renewcommand{\headrulewidth}{0pt} % disable the underline of the header
\raggedbottom
\hfuzz=2mm %suppresses overfull box warnings

\usepackage{hyperref}
\usepackage{float}
\setlength{\parskip}{0.2\baselineskip}
\sisetup{reset-text-series=false, reset-math-version=false, text-series-to-math=true} % \si inherits bold

\title{Physics}
\author{Chris Huson}
\date{October 2026}

\fancyhead[LE]{\thepage}
\fancyhead[RO]{\thepage}
\fancyhead[LO]{La Scuola d'Italia / Huson / Physics \\* 15 October 2026}

\begin{document}

\subsubsection*{2.1 Warm-up: Mental Math Solutions}

\emph{Students write only the answer. The notes give the step a student should be
doing in their head, and the mistakes worth looking for.}

\begin{enumerate}[itemsep=0.45cm]

\item Fill in the blanks.
  \begin{multicols}{2}
    \begin{enumerate}[itemsep=0.2cm]
      \item 2.5 km = \textbf{2500 m}
      \item 450 cm = \textbf{4.5 m}
      \item 0.8 m = \textbf{80 cm}
      \item 35 mm = \textbf{3.5 cm}
      \item 1 m = \textbf{1000 mm}
      \item 7 mm = \textbf{0.007 m}
    \end{enumerate}
  \end{multicols}
  \emph{Going to a smaller unit makes the number bigger. A student who writes
  0.0025 m for (a) or 45{,}000 m for (b) has moved the decimal point the wrong way.}

\item \textbf{10{,}000{,}000 m} (ten million meters).

  \emph{So the distance around the Earth through the poles is about 40{,}000 km. The
  modern value is 40{,}008 km --- the 1790s survey was very good.}

\item \textbf{100 sheets.} \quad 1 cm = 10 mm, and $10 \div 0.1 = 100$.

  \emph{A 500-sheet ream of paper is about 5 cm thick, which students can check.}

\item Phone storage.
  \begin{enumerate}[itemsep=0.2cm]
    \item \textbf{256{,}000 MB.} \quad 1 GB = 1000 MB.
    \item \textbf{64{,}000 photos.} \quad $256{,}000 \div 4 = 64{,}000$.
  \end{enumerate}
  \emph{Phone and drive makers use 1 GB = 1000 MB. Some software uses 1024, so a
  student who knows that may write 262{,}144 MB; accept it, but it is not mental math.}

\item \textbf{4{,}000{,}000{,}000} ticks per second (four billion).

  \emph{Giga means a billion, nine zeros. One hertz is one cycle per second.}

\item Micrometers.
  \begin{enumerate}[itemsep=0.2cm]
    \item \textbf{1000 \si{\micro\metre}} in 1 mm. \quad
      1 m = 1{,}000{,}000 \si{\micro\metre} and 1 m = 1000 mm.
    \item \textbf{125 cells.} \quad $1000 \div 8 = 125$.
    \item \textbf{10{,}000 \si{\micro\metre}.} \quad 1 cm = 10 mm, and
      $10 \times 1000 = 10{,}000$.
  \end{enumerate}
  \emph{About 1250 red blood cells would fit across the chip. The transistors on it
  are smaller again --- hundreds of them would fit across one red blood cell.}

\newpage

\item Fractions as decimals.
  \begin{multicols}{4}
    \begin{enumerate}[itemsep=0.25cm]
      \item $\dfrac{1}{4} = \mathbf{0.25}$
      \item $\dfrac{3}{4} = \mathbf{0.75}$
      \item $\dfrac{1}{5} = \mathbf{0.2}$
      \item $\dfrac{1}{10} = \mathbf{0.1}$
      \item $\dfrac{1}{8} = \mathbf{0.125}$
      \item $\dfrac{5}{8} = \mathbf{0.625}$
      \item $\dfrac{1}{3} = \mathbf{0.333\ldots}$
      \item $\dfrac{2}{3} = \mathbf{0.666\ldots}$
    \end{enumerate}
  \end{multicols}
  \emph{Accept $0.\overline{3}$ and $0.\overline{6}$, and accept 0.33 or 0.67 written
  with $\approx$. Mark ``$\frac{2}{3} = 0.66$'' wrong: it is neither exact nor
  correctly rounded. $\frac{1}{8}$ is half of $\frac{1}{4}$, and
  $\frac{5}{8} = 5 \times 0.125$.}

\item Fractions of an hour and a day.
  \begin{multicols}{2}
    \begin{enumerate}[itemsep=0.2cm]
      \item $\dfrac{1}{4}$ hour = \textbf{15 min}
      \item $\dfrac{1}{3}$ hour = \textbf{20 min}
      \item 0.1 hour = \textbf{6 min}
      \item 30 s = \textbf{0.5 min}
      \item 45 min = \textbf{0.75 h}
      \item 6 hours = $\mathbf{\dfrac{1}{4}}$ \textbf{day}
      \item 8 hours = $\mathbf{\dfrac{1}{3}}$ \textbf{day}
      \item 1.2 hours = \textbf{1 h 12 min}
    \end{enumerate}
  \end{multicols}
  \emph{Part (h) is the one to watch: 1.2 h is not 1 h 20 min, because
  $0.2 \times 60 = 12$. A clock is not a decimal.}

\item \textbf{3600 s} in an hour \quad \textbf{365 days} in a year \quad
  \textbf{52 weeks} in a year

  \emph{$60 \times 60 = 3600$. A year is $365\frac{1}{4}$ days on average, which is
  why we have leap years; accept 365 or 366. $365 \div 7 = 52\frac{1}{7}$, so
  ``about 52.''}

\item Sunlight.
  \begin{enumerate}[itemsep=0.2cm]
    \item \textbf{500 s.} \quad
      $150{,}000{,}000 \div 300{,}000 = 1500 \div 3 = 500$.
    \item \textbf{8 min 20 s.} \quad $8 \times 60 = 480$, and $500 - 480 = 20$.
  \end{enumerate}
  \emph{Cancel the five zeros the two numbers share. In 1-9CW the precise values gave
  8.31 min; these round values give $8\frac{1}{3}$ min.}

\item Circle the best answer.
  \begin{enumerate}[itemsep=0.2cm]
    \item \textbf{(B) 2.54.} \emph{Exact, by definition. The distractors are
      other conversion factors: 1.6 km in a mile, 3.28 ft in a meter, 10 mm in a cm.}
    \item \textbf{(C) $\mathbf{\dfrac{5}{8}}$.} \emph{1 km $\approx$ 0.621 mi, and
      $\frac{5}{8} = 0.625$. Choice (D) is backwards: a mile is about 1.6 km, so a
      kilometer is less than a mile.}
  \end{enumerate}

\item \textbf{About 3 miles.} \quad
  $5 \times \dfrac{5}{8} = \dfrac{25}{8} = 3\dfrac{1}{8}$.

  \emph{Accept 3 or 3.1 mi; the true value is 3.11 mi.}

\end{enumerate}
\end{document}
